\section{A Note on Assign Cookies: When DP Is the Wrong Tool}
\label{sec:dp-assign-cookies-note}

\noindent Striver's sheet lists Assign Cookies again here, inside the DP
on Subsequences lecture -- and it is worth pausing on \emph{why}, since
the honest answer is a valuable lesson in its own right: \textbf{this
problem does not need dynamic programming at all}, and forcing a DP
solution onto it would be strictly worse than the greedy solution
already derived in the Greedy chapter.

\begin{patternbox}[Recognising When a Subset/Assignment Problem Is Actually Greedy]
Assign Cookies asks: given greed factors $g$ and cookie sizes $s$,
maximize the number of children satisfiable, where child $i$ is
satisfied by cookie $j$ if $s[j] \ge g[i]$. This \emph{looks} like an
assignment problem that backtracking or DP might solve via some
subset-matching recurrence -- but the Greedy chapter's Assign Cookies section
already proved, via a direct exchange argument, that sorting both arrays
and greedily matching the smallest sufficient cookie to the smallest
unsatisfied child is \textbf{provably optimal}. No DP table is needed
because there is no genuine trade-off between overlapping sub-problems
to resolve: the greedy choice at each step is never wrong, so there is
nothing for memoization to usefully cache.
\end{patternbox}

\subsection*{Why a DP/backtracking solution would be strictly worse here}

A backtracking solution exploring every possible assignment of cookies to
children would cost at least $O(2^{n})$ in the worst case (each child
either gets a cookie or doesn't, combinatorially), or, with a more
clever subset-based DP state, perhaps $O(n \cdot m)$ or worse depending
on formulation -- but the Greedy chapter's greedy algorithm
already solves the problem in $O(n \log n + m \log m)$ (dominated by
sorting) with $O(1)$ extra space. Since the greedy algorithm is already
proven optimal, any DP formulation can only ever \emph{match} this
bound at best, typically with strictly worse constants and memory
usage -- there is no scenario in which reaching for DP here is the
better engineering decision.

\begin{ideabox}[The Real Lesson]
Before reaching for memoization or tabulation on a subset/assignment-flavored
problem, always ask first: \emph{can a sorted greedy pass already solve
this optimally?} The keyword-recognition habits from the Greedy chapter's
Assign Cookies section, in particular, should be
checked \emph{before} defaulting to a DP recurrence, not after. DP is
the right tool precisely when greedy's exchange argument \emph{fails} --
when the locally best choice can depend on choices made elsewhere, which
is exactly what does \emph{not} happen in Assign Cookies.
\end{ideabox}

\begin{takeawaybox}
Not every problem that involves choosing a subset or making an
assignment needs dynamic programming. Recognizing when a simpler,
already-optimal greedy algorithm applies -- and confidently \emph{not}
reaching for a DP table in that case -- is as important a skill as
knowing how to build the DP table when one is genuinely needed.
\end{takeawaybox}
